% WS4 -- Type C: FRQ practice, rubric visible. Long-FRQ format (this unit's
% exam has a 10-point long FRQ; students should meet the format here first).
\documentclass{apchem-worksheet}

\apcunit{2}
\apcunittitle{Compound Structure and Properties}
\apcdoctitle{Worksheet 4 \textbullet\ FRQ Practice}
\apcsubtitle{Type C \textbullet\ one long FRQ, the exam's format \textbullet\ answer, then self-score}

\begin{document}
\wsheader[The unit exam's FRQ is this shape and length. Write everything
first; the rubric is printed after the last part. Score like a reader: no
point for anything not on the page.]

\problem[2.5]{\textbf{Long FRQ (10 points).} Answer all parts. Nitrogen and
oxygen form several species; this question examines \ce{NO2-} and
\ce{NO2+}.}
\begin{parts}
  \item Give the total valence electron count for each species.
        \answerlines[1]{\ce{NO2-}: $5+12+1=18$. \quad \ce{NO2+}:
        $5+12-1=16$.}
  \item Draw a complete Lewis diagram for \ce{NO2+}. \workspace[2.6cm]
        \keyonly{\begin{apcanswer}O=N=O, linear skeleton, 2 lone pairs per O,
        no lone pair on N; brackets, $+$.\end{apcanswer}}
  \item Draw BOTH resonance forms of \ce{NO2-}. \workspace[2.8cm]
        \keyonly{\begin{apcanswer}O=N--O $\leftrightarrow$ O--N=O, lone pair
        ON N, brackets, $-$.\end{apcanswer}}
  \item Predict the O--N--O bond angle in each species and justify the
        difference. \answerlines[3]{\ce{NO2+}: 180$^\circ$ --- 2 domains,
        linear. \ce{NO2-}: $<120^\circ$ (about 115$^\circ$) --- 3 domains
        including a lone pair on N that compresses the angle.}
  \item In \ce{NO2-}, the two N--O bonds are equal in length. Explain using
        your diagrams from (c). \answerlines[2]{The true structure is the
        resonance hybrid --- an average of the two forms --- so each bond has
        order $3/2$ and identical length.}
  \item Predict which species has the shorter N--O bonds, and justify with
        bond order. \answerlines[2]{\ce{NO2+} --- bond order 2 vs $3/2$ in
        \ce{NO2-}; higher order, shorter bond.}
\end{parts}

\begin{apcnote}[Rubric --- score yourself, 10 points]
\textbf{(a) 1 pt} both counts correct, charge handled both directions.
\textbf{(b) 2 pts} 16 e$^-$ structure: 1 for correct connectivity/multiple
bonds, 1 for lone pairs + charge bracket.
\textbf{(c) 2 pts} both forms, lone pair on N, 18 e$^-$ each.
\textbf{(d) 2 pts} 1 for each angle WITH domain-count justification; bare
angles earn 1 max.
\textbf{(e) 1 pt} hybrid/average language with equal order --- ``flipping
between forms'' earns 0.
\textbf{(f) 2 pts} 1 for \ce{NO2+}, 1 for the order comparison $2 > 3/2$.
\end{apcnote}

\begin{spiralstrip}{Units 1--2 \textbullet\ spiral}
  \item \ce{Mg3N2}: total valence electrons available from 3 Mg atoms:
        \answerblank[2cm]{6}
  \item Rank lattice energy: \ce{LiF}, \ce{KI}, \ce{MgO}:
        \answerblank[4.5cm]{\ce{MgO} $>$ \ce{LiF} $>$ \ce{KI}}
  \item Element with PES peaks 2:2:6:1: \answerblank[2.5cm]{Na}
  \item Moles in \SI{88}{\gram} \ce{CO2}: \answerblank[2.5cm]{\SI{2.0}{\mole}}
  \item Shape of \ce{CO3^2-}: \answerblank[3.5cm]{trigonal planar}
\end{spiralstrip}

\end{document}
